\paragraph{Eigenvalues of Unitary Operators as Phase Factors}\label{ex:eigenvalues-of-unitary-operators-as-phase-factors}

In the previous section, we recalled that an eigenstate \ket{v} of a unitary operator $U$ satisfies:
\begin{equation*}
    U\ket{v}=\lambda\ket{v}
\end{equation*}
Since $U$ is unitary (i.e., $U^\dagger U=I$), the eigenvalue $\lambda$ cannot be an arbitrary complex number. It must have modulus one:
\begin{equation*}
    \left|\lambda\right|=1
\end{equation*}
Therefore, every \textbf{eigenvalue of a unitary operator} can be written as
\begin{equation*}
    \lambda=e^{i\phi}
\end{equation*}
Where $\phi\in\mathbb{R}$ is a phase angle. The eigenvalue equation then becomes:
\begin{equation*}
    U\ket{v}=e^{i\phi}\ket{v}
\end{equation*}
This means that when a unitary operator acts on one of its eigenstates, it \textbf{does not change the direction} of the state vector. It only \textbf{multiplies the state by a complex phase factor}.

\highspace
\begin{flushleft}
    \textcolor{Green3}{\faIcon{question-circle} \textbf{Why must the eigenvalue have modulus one?}}
\end{flushleft}
Suppose that \ket{v} is a normalized eigenstate of $U$:
\begin{equation*}
    U\ket{v}=\lambda\ket{v}
\end{equation*}
Because $U$ is unitary, it preserves the norm of every state. Therefore:
\begin{equation*}
    \left|U\ket{v}\right|=\left|\ket{v}\right|
\end{equation*}
Since \ket{v} is normalized:
\begin{equation*}
    \left\|\ket{v}\right\| = 1
\end{equation*}
Using the eigenvalue equation:
\begin{equation*}
    \left\|U\ket{v}\right\| = \left\|\lambda\ket{v}\right\|
\end{equation*}
Multiplying a vector by a scalar multiplies its norm by the modulus of that scalar:
\begin{equation*}
    \left\|\lambda\ket{v}\right\| = \left|\lambda\right|\left\|\ket{v}\right\|
\end{equation*}
Therefore:
\begin{equation*}
    \left|\lambda\right| \left\|\ket{v}\right\| = \left\|\ket{v}\right\|
\end{equation*}
Since $\left\|\ket{v}\right\|=1$:
\begin{equation*}
    \left|\lambda\right| = 1
\end{equation*}
This is necessary because a quantum gate must preserve normalization and therefore preserve total probability.

\newpage

\begin{flushleft}
    \textcolor{Green3}{\faIcon{compass} \textbf{Complex numbers of modulus one}}
\end{flushleft}
A complex number of modulus one lies on the unit circle of the complex plane. Using Euler's formula:
\begin{equation*}
    e^{i\phi}=\cos\phi+i\sin\phi
\end{equation*}
We obtain:
\begin{equation*}
    \left|e^{i\phi}\right|
    =
    \sqrt{\cos^2(\phi)+\sin^2(\phi)}
    =
    1
\end{equation*}
Therefore, writing:
\begin{equation*}
    \lambda=e^{i\phi}
\end{equation*}
Means that the \textbf{eigenvalue represents a rotation by an angle} $\phi$ in the \textbf{complex plane}, \hl{without changing the magnitude of the state.} The operator changes \hl{only the phase of the eigenstate}:
\begin{equation*}
    U\ket{v}=e^{i\phi}\ket{v}
\end{equation*}
In other words, the \hl{$e^{i\phi}$ changes the phase of the state, but not its probabilities (i.e., the measurement outcomes).}

\highspace
\begin{examplebox}[: Phase eigenvalues]
    Some common eigenvalues are:
    \begin{equation*}
        1=e^{i0},
        \qquad
        -1=e^{i\pi},
        \qquad
        i=e^{i\pi/2},
        \qquad
        -i=e^{-i\pi/2}
    \end{equation*}
    These phase factors correspond to rotations around the unit circle:
    \begin{center}
        \begin{tabular}{@{} c c c @{}}
            \toprule
            Eigenvalue & Exponential form & Phase angle \\
            \midrule
            $1$ & $e^{i0}$ & $0$ \\
            $i$ & $e^{i\pi/2}$ & $\pi/2$ \\
            $-1$ & $e^{i\pi}$ & $\pi$ \\
            $-i$ & $e^{-i\pi/2}$ & $-\pi/2$ \\
            \bottomrule
        \end{tabular}
    \end{center}
    This representation is especially useful because different quantum gates can produce different phase angles on their eigenstates.
\end{examplebox}

\highspace
\begin{examplebox}[: Pauli-$X$]
    The eigenstates of the Pauli-$X$ gate are:
    \begin{equation*}
        \ket{+}=\frac{\ket{0}+\ket{1}}{\sqrt{2}},
        \qquad
        \ket{-}=\frac{\ket{0}-\ket{1}}{\sqrt{2}}
    \end{equation*}
    Their eigenvalue equations are:
    \begin{equation*}
        X\ket{+}=\ket{+}=e^{i0}\ket{+},
        \qquad
        X\ket{-}=-\ket{-}=e^{i\pi}\ket{-}
    \end{equation*}
    Therefore:
    \begin{itemize}
        \item $\ket{+}$ acquires a phase of $0$;
        \item $\ket{-}$ acquires a phase of $\pi$.
    \end{itemize}
    This difference will later be essential in phase kickback. A controlled-$X$ acting on a target in $\ket{-}$ can produce a relative minus sign on the control qubit.
\end{examplebox}

\highspace
\begin{examplebox}[: Phase gate ($S$)]
    The phase gate is:
    \begin{equation*}
        S=
        \begin{pmatrix}
            1&0\\
            0&i
        \end{pmatrix}
    \end{equation*}
    Its computational-basis states are eigenstates:
    \begin{equation*}
        S\ket{0}=\ket{0}=e^{i0}\ket{0},
        \qquad
        S\ket{1}=i\ket{1}=e^{i\pi/2}\ket{1}
    \end{equation*}
    Therefore:
    \begin{itemize}
        \item $\ket{0}$ has eigenvalue $1=e^{i0}$;
        \item $\ket{1}$ has eigenvalue $i=e^{i\pi/2}$.
    \end{itemize}
    The state $\ket{1}$ is not changed into a different basis state. It only acquires a phase of $\pi/2$.
\end{examplebox}

\highspace
\textcolor{Green3}{\faIcon{book} \textbf{Real symmetric unitary operators.}} If a unitary operator is also real and symmetric, then its eigenvalues are real. This simplifies the analysis of eigenvalues, because we already know that every unitary eigenvalue must satisfy:
\begin{equation*}
    \left|\lambda\right|=1
\end{equation*}
The only real numbers with modulus one are:
\begin{equation*}
    \lambda=+1\quad\text{or}\quad\lambda=-1
\end{equation*}
Equivalently:
\begin{equation*}
    +1=e^{i0},
    \qquad
    -1=e^{i\pi}
\end{equation*}
This is why \hl{operators such as $X$, $Z$, and $H$ have eigenvalues $+1$ and $-1$.} The negative eigenvalue should not be interpreted as changing the magnitude of the state. It represents a phase shift of $\pi$:
\begin{equation*}
    -\ket{v}=e^{i\pi}\ket{v}
\end{equation*}
Which does not change the measurement probabilities of the state.

\highspace
In summary, for an eigenstate \ket{v} of a unitary operator $U$:
\begin{equation*}
    U\ket{v}=e^{i\phi}\ket{v}
\end{equation*}
The eigenvalue $e^{i\phi}$:
\begin{itemize}
    \item Has modulus one;
    \item Preserves normalization;
    \item Does not change measurement probabilities by itself;
    \item Represents a phase rotation;
    \item Can become computationally useful when converted into a relative phase.
\end{itemize}
